% Generated draft: requires mathematical and source review.
\documentclass[11pt]{article}
\usepackage[margin=25mm]{geometry}
\usepackage{fontspec}
\setmainfont{Latin Modern Roman}
\usepackage{amsmath,amssymb,mathtools,unicode-math}
\setmathfont{Latin Modern Math}
\usepackage{enumitem,needspace,xurl,hyperref}
\setlength{\parindent}{0pt}
\setlength{\parskip}{5pt}
\setlength{\emergencystretch}{3em}
\newcommand{\N}{\mathbb{N}}
\newcommand{\Z}{\mathbb{Z}}
\newcommand{\Q}{\mathbb{Q}}
\newcommand{\R}{\mathbb{R}}
\newcommand{\C}{\mathbb{C}}
\newcommand{\F}{\mathbb{F}}
\newcommand{\A}{\mathbb{A}}
\newcommand{\PP}{\mathbb{P}}
\newcommand{\E}{\mathbb{E}}
\newcommand{\eps}{\varepsilon}
\newcommand{\dd}{\,\mathrm{d}}
\begin{document}
\section*{UnsolvedMath ID: 20000741}
\textbf{Source problem code:} catalogue entry
\subsection{Definitions and mathematical statement}
This permanent record splices two distinct workshop items and must not be treated as one theorem. The supplied verification identifies the apparent heading ``14'' as footer page number \(1\) followed by Problem \(4\). Both mathematical directions are retained under the one permanent ID without inventing a relation between them.

\textbf{Problem 3: restricted-modulus large sieve.} Y. Motohashi proposes improving a large-sieve inequality by restricting its moduli and mentions D. Wolke's results for prime moduli. The source gives no modulus class, weights, or desired gain. One standard model to be specified is an estimate for
\[
\sum_{\substack{q\leq Q\\q\in\mathcal Q}}
\sum_{\substack{1\leq a\leq q\\(a,q)=1}}
\left|\sum_{n=M+1}^{M+N}c_ne^{2\pi i an/q}\right|^2,
\]
relative to \(\sum|c_n|^2\), improving the usual order \(N+Q^2\) in a suitable range by using information about \(\mathcal Q\). This is an illustrative additive large-sieve formulation, not a uniquely recovered quantitative target; the original does not distinguish it from other large-sieve variants or authorize a universal gain for every sparse set.

\textbf{Problem 4: the zeta inequality.} For real \(t\ne0\), put \(F(t)=|t\zeta(1+it)|\), and use the continuous extension \(F(0)=1\) from the simple pole of \(\zeta\) at \(1\). J. Pintz asks whether
\[
F(t)\geq1\qquad(t\in\R).
\]
The source notes the inequality for sufficiently small \(|t|\) and sufficiently large \(|t|\), and says a complete proof would simplify some GPY arguments. R. C. Vaughan recalls an affirmative solution in his files; this recollection is not supplied as a recoverable written proof.

The attached aid reports that C. Y. Yıldırım's 2026 theorem proves \(F(t)>1\) for every nonzero real \(t\), without giving its full citation. Its separate local calculation uses the Stieltjes convention
\[
\zeta(1+z)=z^{-1}+\sum_{j=0}^\infty\frac{(-1)^j\gamma_j}{j!}z^j,
\qquad\gamma_0=\gamma,
\]
and reports
\[
F(t)^2=1+(\gamma^2+2\gamma_1)t^2
+(\gamma_1^2-\gamma\gamma_2-\gamma_3/3)t^4+O(t^6).
\]
It states that rational bounds \(\gamma>1/2\) and \(\gamma_1>-1/8\) certify \(\gamma^2+2\gamma_1>0\), so the extended function has a strict local minimum at \(0\). It also reports the explicit sufficient range \(|t|\geq250\) from a published reciprocal-zeta estimate, whose exact statement and reference are absent. These supplied unreviewed claims address Problem 4 only. They neither make the underspecified Problem 3 precise nor prove a restricted-modulus large-sieve improvement.
\subsection{Short English statement}
This record combines an underspecified large-sieve direction with Pintz's separate inequality \(|t\zeta(1+it)|\geq1\). The latter needs a continuous interpretation at zero and is reported solved by a 2026 theorem; the local expansion and large-\(|t|\) bound do not address the other item.
\subsection{Sources}
\begin{itemize}
\item AIM workshop \emph{Gaps between primes}, consecutive Problems 3 and 4: \url{https://aimath.org/WWN/primegaps/WorkshopProblems.pdf}. The source attributes the restricted-modulus direction to Y. Motohashi, mentions D. Wolke, and attributes the zeta question and historical comments to J. Pintz and R. C. Vaughan.
\item Local registry \texttt{data/unsolvedmath/problems.json}, permanent ID \(20000741\), exact source code \texttt{catalogue entry}; canonical \texttt{aim-analytic-number-theory-notes.json}, notes[104].
\item The attached aid names C. Y. Yıldırım's 2026 theorem and a published reciprocal-zeta bound without full bibliography. Attempt 1 is explicitly machine-generated, not peer reviewed, and requires independent verification before citation.
\end{itemize}
\end{document}
